\documentclass[profile=standard]{ipara-handbook}
\title{B06A｜CDF 与 PDF：局部概率密度与累积}
\author{}
\date{}
\begin{document}
\maketitle

\begin{quote}
本册回答：\textbf{为什么概率论里的 CDF 与 PDF 会呈现“积分—求导”关系？为什么这条关系对连续密度非常自然，却不能把所有分布都粗暴地写成 $F'=f$？}
\end{quote}

\section{本册定位：把概率总账与局部强度接到微积分基本定理}
概率 Topic03 Own 随机变量、分布、CDF、PMF 与 PDF；高数 Topic05/H-B03 Own 微分—累积的基本定理与正则性边界。B06A 只 Own 一条跨科接口：
\[
\boxed{
\text{概率质量在 }(-\infty,x]\text{ 上的累积}
\xrightarrow{\text{CDF}}
F(x)
\quad\Longleftrightarrow\quad
\text{局部连续质量强度 }f(x)
}
\]

在具有密度的情形，二者通过
\[
F(x)=\int_{-\infty}^{x}f(t)\,dt
\]
连接；在适当正则点，微积分基本定理给
\[
F'(x)=f(x).
\]

\section{先分三层对象：分布、CDF、PDF}
分布本身是概率质量怎样落在实轴上的整体规则。CDF 是它的通用累计表示：
\[
F_X(x)=P(X\le x).
\]
任何实随机变量都有 CDF。

PDF 则只属于具有绝对连续分布的情形：若存在非负可积函数 $f$ 使
\[
P(X\in A)=\int_A f(x)\,dx
\]
对相应集合成立，则 $f$ 是密度。

因此对象关系应明确写成：\textbf{每个实随机变量的分布都可以用 CDF 表示；只有具有绝对连续分布的那一类，才进一步拥有普通 PDF 表示。}不能把“每个分布都有 CDF”错误升级成“每个 CDF 都有普通导数密度”。

\section{为什么密度值不是概率}
对很短区间 $[x,x+\Delta x]$，若 $f$ 在 $x$ 附近连续，则
\[
P(x<X\le x+\Delta x)
\approx f(x)\Delta x.
\]
所以 $f(x)$ 是“单位长度上的局部质量强度”，真正概率必须再乘一个长度并累积。

因此密度可以大于 $1$，只要总积分仍为 $1$；而连续型随机变量通常满足
\[
P(X=x)=0.
\]

这也是为什么
\[
\boxed{f(x)\neq P(X=x).}
\]

\section{微积分基本定理怎样进入概率}
若
\[
F(x)=\int_{-\infty}^{x}f(t)\,dt,
\]
则对 $a<b$，
\[
P(a<X\le b)=F(b)-F(a)=\int_a^b f(t)\,dt.
\]

当 $f$ 在 $x$ 处连续时，FTC 给出
\[
F'(x)=f(x).
\]
更一般地，若 $F$ 绝对连续，则存在 $f\in L^1$ 使
\[
F(x)=F(-\infty)+\int_{-\infty}^{x}f(t)dt
\]
并且 $F'=f$ 几乎处处。

对数学一来说不需要展开完整测度论，但这个背景非常重要：\textbf{“PDF 是 CDF 的导数”其实依赖绝对连续结构，不是 CDF 的定义。}

\section{跳跃为什么代表原子质量}
若 CDF 在 $x_0$ 处跳跃，则
\[
P(X=x_0)=F(x_0)-F(x_0^-).
\]
这部分质量集中在单点上，无法由普通函数密度在零长度单点上的积分产生。

所以离散分布与混合分布必须保留“跳跃账本”。例如
\[
P(X=0)=\frac12,
\]
其余一半在 $(0,1)$ 上均匀，则 CDF 同时包含跳跃与连续增长。对连续段求导只能得到连续部分密度，不能把原子质量弄丢。

\section{母例：指数分布把局部—累积完整跑一遍}
若
\[
f(x)=e^{-x}\mathbf 1_{x\ge0},
\]
则
\[
F(x)=
\begin{cases}
0,&x<0,\\
1-e^{-x},&x\ge0.
\end{cases}
\]
对 $x>0$，
\[
F'(x)=e^{-x}=f(x).
\]
区间概率既可由 CDF 做差，也可由密度积分：
\[
P(a<X\le b)=F(b)-F(a)=\int_a^b e^{-x}dx.
\]

这个母例的重点不是积分本身，而是看清：CDF 是累计账本，PDF 是局部连续强度，FTC 负责两者之间的合法转换。

\section{B06A 串起哪些章节}
\begin{tabular}{p{0.23\linewidth}p{0.30\linewidth}p{0.39\linewidth}}
\hline
Owner / 下游 & 提供对象 & B06A 的翻译责任\\
\hline
高数 Topic05 / H-B03 & 变上限积分、FTC、正则性 & 解释“局部强度积成累计量、累计量何时可微”\\
概率 Topic03 & CDF、PMF、PDF、混合分布 & 解释哪些概率对象能走积分—求导接口，哪些必须保留跳跃\\
B06B & 概率积分语言 & 向联合/边缘/期望输出“质量必须通过积分/求和汇总”的基础语义\\
B07 & 变量变换 & 提供新分布最后仍须通过 CDF/密度账本验收的接口\\
\hline
\end{tabular}

\section{反桥梁：四个最危险的越级}
\begin{tabular}{p{0.20\linewidth}p{0.04\linewidth}p{0.20\linewidth}p{0.48\linewidth}}
\hline
A & & B & 真正区别\\
\hline
CDF 连续 & $\neq$ & 存在普通 PDF & 连续只能排除跳跃，仍可能存在奇异连续分布\\
PDF 值 & $\neq$ & 点概率 & 密度是局部强度，概率是对区间/集合的积分\\
$F'=f$ & $\neq$ & 所有点处必成立 & 一般绝对连续情形只需几乎处处；不连续密度点要单独审计\\
零点概率 & $\neq$ & 不可能取到 & 连续变量单点概率为零，但该点仍可在支撑中\\
\hline
\end{tabular}

\section{调用信号与退出边界}
看到“由 $F$ 求 $f$、由 $f$ 求 $F$、区间概率、分布函数合法性、混合分布”时，先问：
\begin{enumerate}
\item 当前对象是 CDF、PDF 还是点质量？
\item CDF 是否有跳跃？连续部分是否具备密度资格？
\item 当前目标用做差最直接，还是需要积分/求导？
\item 最后总质量是否为 $1$，CDF 是否单调、右连续并满足两端极限？
\end{enumerate}

\section{一页压缩}
B06A 最终只保留三句：
\[
\boxed{F(x)=P(X\le x)}
\]
是总账；若分布具有密度，则
\[
\boxed{F(x)=\int_{-\infty}^{x}f(t)dt};
\]
在合适正则点再有
\[
\boxed{F'(x)=f(x)}.
\]

真正要记住的边界是：\textbf{跳跃质量不在普通密度里，CDF 总存在而 PDF 不总存在。}

\end{document}
